% problem-000150 9 有机合成与机理
\section*{9. 有机合成与机理}
瑞德西⻙(remdesivir)是吉利德公司开发的⼀种新型核酸类抗病毒药物。2021年11⽉，默克公司研发的抗病毒药物莫努匹⻙(molnupiravir)在英国上市。Remdesivir 和 molnupiravir 的抗病毒作⽤机制都是通过阻断RNA聚合酶(RdRp)的合成抑制病毒的复制。

（图：）

近期，我国科学家对remdesivir的合成路线进⾏了重要改进：通过发展⼀种⼿性咪哇类有机⼩分⼦催化剂，解决了remdesivir的磷⼿性中⼼⾼效构建的难题。其中部分合成路线如下：

（图：）

\subsection*{9-1}
标出remdesivir分⼦结构中1号位碳(C1)和磷的绝对构型。

\subsection*{9-2}
画出第⼀步反应中间体4（未经分离提纯，直接⽤于下⼀步反应）和原料5的结构简式，并解释第⼀步反应的作⽤。

\subsection*{9-3}
本路线所合成的化合物6是⼀对C1位差向异构体的混合物 $( \beta \mathrm { { : a } } = 3 \mathrm { { : 1 } ) }$ ，但化合物7为单⼀构型产物。画出的结构简式，并利⽤反应中间体解释6转化为单⼀构型7的历程。

\subsection*{9-4}
Molnupiravir存在两种互变异构体，可分别模仿胞苷C和尿苷U，导致病毒RNA复制过程的突变率增加。写出 molnupiravir 的另⼀种互变异构体，判断这两种互变异构体分别参与的主要碱基配对⽅式，并画出其⽰意图。

（图：）
